\documentclass[tikz,border=4pt]{standalone}
% Paleta y fuentes compartidas por el libro y las figuras (única fuente de verdad).
\RequirePackage{xcolor}
\definecolor{oxido}{HTML}{B7410E}
\definecolor{cobre}{HTML}{D27D2D}
\definecolor{carbon}{HTML}{2B2B2B}
\definecolor{azul}{HTML}{1F5F8B}
\definecolor{arena}{HTML}{F4EDE4}

\usepackage{fontspec}
\setsansfont{FiraSans}[Extension=.otf,UprightFont=*-Regular,BoldFont=*-SemiBold,ItalicFont=*-Italic,BoldItalicFont=*-SemiBoldItalic]
\setmonofont{FiraMono}[Extension=.otf,UprightFont=*-Regular,BoldFont=*-Medium,Scale=0.95]
\usetikzlibrary{arrows.meta,calc,positioning,shapes.geometric}
\begin{document}
\sffamily\small
% Retos del informe Vision Doc (marzo de 2026) y trabajo asignado en 2026.
\begin{tikzpicture}[
  reto/.style={draw=#1,line width=0.9pt,fill=white,rounded corners=3pt,
    minimum height=16mm,text width=50mm,align=left,inner sep=6pt,
    execute at begin node={\hyphenpenalty=10000\raggedright}},
  plan/.style={fill=oxido,text=white,rounded corners=2pt,inner sep=2.5pt,
    font=\sffamily\scriptsize\bfseries},
  nota/.style={font=\sffamily\footnotesize,text=carbon!75}]
  \node[reto=oxido] (r1) at (0,0) {{\bfseries\color{oxido}Tiempos de compilación}\\[2pt]
    {\footnotesize Objetivos: compilador paralelo y nuevo resolvedor de rasgos.}};
  \node[reto=oxido] (r2) at (6.2,0) {{\bfseries\color{oxido}Programación asíncrona}\\[2pt]
    {\footnotesize Hoja de ruta «Just add async».}};
  \node[reto=carbon!55] (r3) at (0,-2.2) {{\bfseries\color{carbon}Propiedad y préstamos}\\[2pt]
    {\footnotesize La curva de entrada para quien empieza.}};
  \node[reto=carbon!55] (r4) at (6.2,-2.2) {{\bfseries\color{carbon}Elección de \emph{crates}}\\[2pt]
    {\footnotesize Saber qué biblioteca usar y en cuál confiar.}};
  \node[plan] at (r1.north east) [anchor=center,xshift=-9mm] {objetivo 2026};
  \node[plan] at (r2.north east) [anchor=center,xshift=-9mm] {objetivo 2026};
  % Dominios inmaduros, en una franja inferior.
  \node[draw=azul,dashed,line width=0.9pt,rounded corners=3pt,fill=azul!6,
    text width=112mm,align=center,inner sep=6pt,
    execute at begin node={\hyphenpenalty=10000}] (dom) at (3.1,-4.25)
    {{\bfseries\color{azul}Dominios inmaduros}\\[2pt]
     {\footnotesize embebidos con pocos recursos, seguridad crítica e interfaces gráficas (capítulo~III)}\\[2pt]{\footnotesize\color{carbon!75}La seguridad crítica tiene hoja de ruta propia; las interfaces gráficas, ningún plan visible.}};
  \node[nota,anchor=south west] at ($(r1.north west)+(0,3mm)$)
    {Cuatro retos comunes, según el informe del \emph{Vision Doc} (marzo de 2026)};
\end{tikzpicture}
\end{document}
